\lecture{Formulario}{26~-~01~-~2026}
\section*{Formulario}

\subsection*{Progettazione e Dimensionamento Plant Layout}
\begin{itemize}
    \item \textbf{Takt Time (Base):}
          \[
              TT = \frac{\text{Tempo Disponibile}}{\text{Domanda}} \quad \left[\frac{\text{tempo}}{\text{unità}}\right]
          \]

    \item \textbf{Frequenza di produzione:}
          \[
              f = \frac{1}{TT} \quad \left[\frac{\text{unità}}{\text{tempo}}\right]
          \]

    \item \textbf{Tempo di lavoro effettivo:}
          \[
              T_{eff} = T_{disp} \times \text{Rendimento}
          \]

    \item \textbf{Unità da produrre (con scarti):}
          \[
              u = \frac{\text{Numero di unità}}{(1-r_{0})(1-r_{1}) \dots (1-r_{n})} \quad \text{(dove } r_n = \text{tasso scarto)}
          \]

    \item \textbf{Pitch Time:}
          \[
              PT = TT \times Q_{contenitore}
          \]

    \item \textbf{Numero Macchine Teorico:}
          \[
              \#{\text{mac.}} =\frac{\text{Tempo Ciclo Macchina}}{\text{Takt Time}} \quad \text{(Arrotondare per eccesso)}
          \]

    \item \textbf{Funzione Obiettivo (Costo Trasporti):}
          \[
              C = \sum_{i} \sum_{j} T_{ij} \cdot c_{ij} \cdot d_{ij}
          \]

    \item \textbf{Rapporto di Hollier (Sequenziamento):}
          \[
              r = \frac{\text{From (Uscite)}}{\text{To (Entrate)}}
          \]
\end{itemize}

\subsection*{Indici di Prestazione (KPI) di Magazzino}
\begin{itemize}
    \item \textbf{Selettività:}
          \[
              I_{s} = \frac{\text{Codici Dir. Accessibili}}{\text{Codici Totali}} \quad [\%]
          \]

    \item \textbf{Coefficiente di Utilizzazione Superficiale:}
          \[
              C_{s} = \frac{S_{\text{stoccaggio}}}{S_{totale}} \quad [\%]
          \]

    \item \textbf{Coefficiente di Utilizzazione Volumetrico:}
          \[
              C_{v} = \frac{V_{\text{stoccaggio}}}{V_{totale}} \quad [\%]
          \]

    \item \textbf{Giacenza Media:}
          \[
              G_{\text{media}} = \frac{\sum q_{i}t_{i}}{\sum t_{i}} \quad [\text{UdC} \text{ o } m^3]
          \]

    \item \textbf{Indice di Rotazione:}
          \[
              I_{R} = \frac{U_{\text{transito}}}{G_{\text{media}}}
          \]

    \item \textbf{Indice di Durata:}
          \[
              I_{D} = \frac{1}{I_{R}} \quad [\text{Periodo T}]
          \]

    \item \textbf{Giorni di Copertura:}
          \[
              GC = \frac{GG_{\text{lavorativi}}}{I_{R}} \quad [\text{Giorni}]
          \]

    \item \textbf{Indice di Honeycombing:}
          \[
              I_{H} = \frac{\text{Ubicazioni Vuote}}{\text{Ubicazioni Totali}} \quad [\%]
          \]
\end{itemize}

\subsection*{Dimensionamento Sistema di Movimentazione (Carrelli)}
\begin{itemize}
    \item \textbf{Tempo Ciclo Totale:}
          \[
              T_{ciclo} = T_{fisso} + T_{variabile}
          \]

    \item \textbf{Percorso totale medio ($r$):}
          \[
              r = n \cdot (2w+a)+ 2 \cdot (2A+m \cdot l)
          \]

    \item \textbf{Durata del ciclo ($d$):}
          \[
              d = \frac{r}{v_{t}} + \frac{2 \cdot h \cdot (j-1)}{v_{s}} + t_{f}
          \]

    \item \textbf{Numero di carrelli necessari ($z$):}
          \[
              z = \frac{d \cdot f}{g \cdot y \cdot \eta}
          \]
\end{itemize}

\subsection*{Sistemi di Stoccaggio e Picking (AS/RS)}
\begin{itemize}
    \item \textbf{Tempi di percorrenza (Orizzontale/Verticale):}
          \[
              T_l = \frac{L}{v_l}, \quad T_h = \frac{H}{v_h}
          \]

    \item \textbf{Fattore di forma ideale:}
          \[
              h_l = \frac{T_h}{T_l} = 1
          \]

    \item \textbf{Punti Caratteristici (FEM 9851):}
          \[
              P_1 = \left(\frac{1}{5}L, \frac{2}{3}H\right), \quad P_2 = \left(\frac{2}{3}L, \frac{1}{5}H\right)
          \]

    \item \textbf{Tempo Medio Ciclo Semplice (CS):}
          \[
              T_m(CS) = t(P_1) + t(P_2) + 2t_{fissi}
          \]
          \emph{Nota: $t(P) = \max(t_x, t_y)$ (movimenti simultanei)}

    \item \textbf{Tempo Medio Ciclo Combinato (CC):}
          \[
              T_m(CC) = T_m(CS) + t(P_1 \to P_2) + 2t_{fissi}
          \]

    \item \textbf{Numero Trasloelevatori:}
          \[
              N = \frac{PM}{\text{Cicli}_{eff}} \quad (\text{dove } PM = \text{Potenzialità richiesta})
          \]

    \item \textbf{Indice di Accesso (Picking):}
          \[
              I_A = \frac{\text{Popularity}}{\text{Volume}} \quad \left[\frac{\text{accessi}}{m^3}\right]
          \]
\end{itemize}

%
\subsection*{Illuminazione}
\begin{itemize}
    \item \textbf{Metodo del flusso totale:}
          \[
              F_{t} = \frac{E \cdot A}{u \cdot m} \quad [\si{lumen}]
          \]
          Dove: $ F_{t} $ è flusso totale delle lampade installate, $ E $ è illuminamento medio previsto ($ \si{lux} $), $ A $ area del locale ($ \si{m^2} $), $ u $ coefficiente di illuminazione ($ < 1 $), $ m $ fattore di manutenzione ($ < 1 $).
    \item \textbf{Numeri di apparecchi di illuminazione previsti:}
          \[
              N = \frac{F_{t}}{F_{L}} = \frac{E \cdot A}{u \cdot m \cdot F_{L}}
          \]
          Dove: $ N $ è il numero totale di apparecchi di illuminazione previsti, $ F_{t} $ il flusso totale delle lampadine installate ($ \si{lumen} $) e $ F_{L} $ il flusso iniziale emesso da un apparecchio di illuminazione ($ \si{lumen} $).
\end{itemize}
